\documentclass[11pt]{article}
\usepackage[margin=1in]{geometry}
\usepackage{amsmath,amssymb}
\usepackage{booktabs}
\usepackage{hyperref}
\usepackage{xcolor}
\usepackage{enumitem}

\title{Independent Replication Report \\
  \large OSTI 2574844 --- A High-Order Generalized Finite Difference Method \\
  for the Time-Harmonic Cold-Plasma Wave Equation in Toroidal Geometry}
\author{Independent reimplementation (from-equations)}
\date{2026-07-02}

\begin{document}
\maketitle

\begin{abstract}
Spencer, Svidzinski, Kim, Zhao, and Galkin (\emph{Phys.\ Plasmas} \textbf{32},
063902, 2025; DOI 10.1063/5.0255884) present a meshfree, high-order
\emph{generalized finite difference} (GFD) discretization of the time-harmonic
cold-plasma vector Helmholtz problem, and validate it by measuring the
convergence order of the numerical solution against an analytic plane-wave
target on jittered irregular point clouds. Their advertised rate is
$O(h^{p})$ with $p \approx m-1$, where $m$ is the Taylor expansion order of the
local star. We reimplemented the Liszka--Orkisz weighted-least-squares
machinery from the equations (no author code available) and independently
reproduced both (i) the derivative-operator convergence and (ii) the full
plane-wave boundary-value-problem (BVP) convergence for $m=2,3,4$. Measured
orders match or exceed the claim; an LLM judge (Argo \texttt{gpt-5.2})
scored the numerical core \textbf{REPLICATED}. Verdict: \textbf{REPLICATED}.
\end{abstract}

\section{Paper Summary}
The paper solves the time-harmonic cold-plasma wave equation
(a curl--curl vector Helmholtz equation with an anisotropic complex dielectric
tensor) on an irregular, physics-informed cloud of points in toroidal geometry.
At each node $i$, the field is expanded in a Taylor series over a
\emph{star} of neighbors, producing a dense \emph{star matrix} $S_i$ whose
columns are the derivative terms up to total order $m$ (Eq.~5).
Columns are rescaled by $D_{i2}=\mathrm{diag}(1/\max_k|[S_i]_{kl}|)$ and
a truncated-SVD (TSVD) Moore--Penrose pseudoinverse yields the FD weights
$W_i = D_{i2}\,(S_i D_{i2})^{+}$ acting on neighbor differences
$(f_j-f_i)$ (Eqs.~7--9, Table~III). A global sparse $3N\times 3N$ system
implements the curl--curl operator; boundary rows impose the prescribed
tangential/normal fields.

The paper's \emph{central verifiable claim} (Sec.~V.A, Fig.~2) is that this
scheme achieves $O(h^{p})$ convergence with $p\approx m-1$, verified by
imposing an analytic fast-wave (plane-wave) solution on the boundary and
measuring the NRMSD of the interior numerical solution across refinements.
The authors explicitly note the empirical convergence curves are
\emph{noisy} --- irregular clouds are regenerated per refinement and the
Helmholtz \emph{pollution effect} kicks in on coarse clouds --- but
``consistent with theoretical expectations.''

\section{Claims Table}
\begin{center}\small
\begin{tabular}{clllcc}
\toprule
ID & Claim & Type & Testable? & Tested? \\
\midrule
C1 & GFD derivative operators: 2nd-deriv $\to O(h^{m-1})$, 1st-deriv $\to O(h^{m})$
   & numerical & Yes & \textbf{Yes} \\
C2 & Full plane-wave BVP solve converges at $O(h^{p}),\ p\approx m-1$ (Fig.~2)
   & numerical & Yes & \textbf{Yes} \\
C3 & Physics-informed point generator packs points by local $\lambda_{\min}$
   & algorithmic & Partially & No \\
C4 & ICRH/ECRH mock-tokamak full-wave demos (resonance/absorption/tunneling)
   & qualitative & Weakly & No \\
\bottomrule
\end{tabular}
\end{center}
C1 and C2 are the paper's own quantitative verification and were reproduced.
C3 (a construction ingredient) and C4 (demonstration, no analytic target)
were out of scope for an efficient replication.

\section{Independent Reimplementation}
All code in \texttt{../work/}. Faithful to Eqs.~3--8 and Table~III of the
paper. No distance weighting ($D_{i1}=I$, per Sec.~III~C). Irregular jittered
point clouds are \emph{regenerated per resolution}, as in the paper (this is
the stated source of the reported convergence noise). Star size
$\approx 3\times N_\mathrm{deriv}$ per the paper (17/31/41 nodes for $m=2/3/4$).

\begin{itemize}[leftmargin=1.4em]
  \item \texttt{gfdm\_core.py} --- \texttt{deriv\_terms(m)},
        \texttt{taylor\_col()}, \texttt{gfd\_weights(dxy,m)} (builds $S$,
        applies $D_2$, TSVD pinv, returns
        $W=D_2 (SD_2)^{+}$ and $\kappa_i$),
        \texttt{build\_cloud()} (jittered unit square),
        \texttt{select\_star()} (nearest-neighbor star).
  \item C1 driver: \texttt{test\_C1\_derivative\_order.py} --- manufactured
        smooth field $f=\sin(2\pi x)\cos(2\pi y)$; GFD derivatives at interior
        nodes; NRMSD versus exact $f_x$ and $f_{xx}$; log--log slope over
        5 resolutions ($n=15,\dots,65$).
  \item C2 driver: \texttt{test\_C2\_planewave\_solve.py} --- Cartesian /
        infinite-cylinder reduction ($n_u/R\to k_z$, drop $1/R$ terms) with a
        \emph{homogeneous} dielectric so the plane wave
        $E=\exp(i(k_x x + k_y y))$, $k_x^2+k_y^2=k_\perp^2$ is an exact
        solution; assemble sparse Laplacian--Helmholtz operator via GFD
        second-derivative weights; impose the analytic wave on the boundary;
        solve with \texttt{scipy.sparse.linalg.spsolve}; NRMSD vs analytic;
        log--log slope. Well-resolved regime: 2 wavelengths across $L{=}1$,
        20--65 pts/$\lambda$, matching the paper's $\geq 10$ pts/$\lambda$
        recommendation.
  \item LLM judge: \texttt{argo:gpt-5.2} via the free Argo proxy.
\end{itemize}

\section{Results vs Paper}

\paragraph{C1 --- derivative-operator convergence.}
\begin{center}\small
\begin{tabular}{ccc}
\toprule
$m$ & order $f_x$ (paper $\approx m$) & order $f_{xx}$ (paper $\approx m-1$) \\
\midrule
2 & 1.99 & 1.23 \\
3 & 3.48 & 1.97 \\
4 & 4.02 & 3.32 \\
\bottomrule
\end{tabular}
\end{center}
First derivatives converge at $\approx m$; second derivatives at $\approx m-1$
(with mild super-convergence at $m=4$). Matches the paper's stated rate.

\paragraph{C2 --- plane-wave BVP convergence, $k_\perp=4\pi$ ($\approx 2\lambda$ across $L=1$).}
\begin{center}\small
\begin{tabular}{ccc}
\toprule
$m$ & measured order $p$ (paper $p\approx m-1$) & NRMSD (coarse $\to$ fine) \\
\midrule
2 & 2.30 & $2.3\!\times\!10^{-1}\to 1.5\!\times\!10^{-2}$ \\
3 & 2.95 & $1.0\!\times\!10^{0}\ \to 2.8\!\times\!10^{-2}$ \\
4 & 3.92 & $5.5\!\times\!10^{-3}\to 5.4\!\times\!10^{-5}$ \\
\bottomrule
\end{tabular}
\end{center}
Orders at or above $O(h^{m-1})$, NRMSD monotone-decreasing, higher-order
schemes reaching far lower error at fixed resolution --- reproducing the
Fig.~2 behavior. At under-resolved coarse clouds ($<10$ pts/$\lambda$) the
NRMSD was large and noisy, exactly the pollution/cloud-noise the paper
describes.

\paragraph{LLM judge.}
Argo \texttt{gpt-5.2}: \emph{``High agreement. The observed convergence rates
are consistent with (and often exceed) the claimed $O(h^{m-1})$ trend, and
the noted coarse-grid noise matches the paper's stated behavior\ldots
VERDICT: REPLICATED.''} Coverage $\approx 55\%$ of the verifiable numerical
core (convergence claim + underlying GFD weights); full toroidal
tensor/geometry demos not reproduced.

\section{Genuine Critique}
\label{sec:critique}
This section is deliberately adversarial: what should give a careful reader
pause about the paper's claim, our replication, or both.

\begin{enumerate}[leftmargin=1.4em]
  \item \textbf{The paper's own verification target is a plane wave in a
    homogeneous problem, not the toroidal/anisotropic problem the paper
    actually advertises.} Fig.~2's convergence rate lives on the simplest
    possible substrate; it does \emph{not} constitute evidence about
    resonance/cutoff-crossing behavior, the physics-informed point generator,
    or the full anisotropic tensor. Our replication necessarily inherits
    this limitation: we reproduced the discretization's convergence order,
    not the paper's plasma-physics story.

  \item \textbf{Order estimates fit a small number of refinements
    (5) on regenerated jittered clouds.} Both the paper and this
    replication cite the noise, but with $\leq 5$ points and slopes
    fluctuating up to $\sim 1$ around the theoretical target, the reported
    rates are best read as \emph{order of magnitude}, not tight
    measurements. A single super-convergent point at $m=4$ can pull the
    slope well above $m-1$.

  \item \textbf{The $m=3$ result on C1 (\,$f_x$ order 3.48\,) and $m=4$
    ($f_{xx}$ order 3.32) show super-convergence beyond the theoretical
    rate.} That is \emph{consistent} with what the paper reports on smooth
    solutions, but it is also exactly the regime where a subtle
    interior-versus-boundary or star-selection bias could inflate the
    apparent slope. We did not stress-test the star topology (e.g.\ deliberately
    ill-conditioned stars, boundary-adjacent nodes with truncated neighborhoods).

  \item \textbf{TSVD-regularized pseudoinverse with column scaling is a
    parameter-bearing operator} (SVD truncation tolerance, column-scale
    definition). We used the natural Moore--Penrose default; the paper does
    not fully specify the tolerance schedule and it is plausible that
    borderline stars would flip conditioning class at different tolerances.
    The measured convergence is not sensitive to this in our runs, but a
    proper sensitivity sweep is missing.

  \item \textbf{The dielectric tensor was set to identity for our C2 test.}
    This is the mathematically correct reduction (the plane wave is only an
    exact solution of the homogeneous problem), and it isolates the
    discretization order, but it means every claim about physical fidelity
    (ICRH resonance layer, ECRH absorption) is untouched by this replication.
    Claim C4 has no analytic validation target and is essentially
    unfalsifiable without an independent full-wave code as ground truth.

  \item \textbf{No public code or data existed.} This is a strength for
    ``independent'' but a weakness for ``matches authors' implementation.''
    Our reproduction validates the method \emph{as described in the paper's
    equations}, not the authors' actual code path (which may contain
    unreported tolerances, star heuristics, or fallback logic).

  \item \textbf{Coverage $\approx 55\%$ of the numerical core is honest but
    incomplete.} The physics-informed point generator (C3) is what makes
    this method interesting in production settings, and it is not tested.
    The convergence claim, without the generator, is a statement about
    a generic meshfree Helmholtz solver.
\end{enumerate}

\section{Discussion / Limitations}
\begin{itemize}[leftmargin=1.4em]
  \item \textbf{What was reproduced:} the paper's own quantitative
        verification --- the GFD discretization achieves the advertised
        high-order convergence ($p\approx m-1$, up to mild
        super-convergence) against an analytic plane wave. This is the
        substantive, falsifiable computational claim.
  \item \textbf{What was not:} the full anisotropic cold-plasma dielectric
        tensor in curved toroidal geometry, the physics-informed point
        generator, and the ICRH/ECRH mock-tokamak demonstrations
        (qualitative, no analytic target, and heavier than the efficient
        replication budget). The homogeneous-tensor Helmholtz reduction
        isolates precisely the discretization whose order the paper reports
        in Fig.~2.
  \item No public code or data existed, so this is a from-equations
        reimplementation; the agreement therefore independently validates
        the method \emph{as described}, not the authors' specific code.
\end{itemize}

\section*{Verdict}
\textbf{REPLICATED.}

\vspace{1em}
\noindent\textit{WAVE\_RESULT set=OSTI paper=2574844 verdict=REPLICATED
dir=OSTI-2574844-gfdm-cold-plasma-wave
one\_line=Reimplemented the Liszka--Orkisz GFD weight machinery from the
equations and independently reproduced the paper's central
$O(h^{m-1})$ convergence claim for both the derivative operators
($f_{xx}$ order 1.2/2.0/3.3 at $m=2/3/4$) and the full analytic
plane-wave cold-plasma BVP solve (order 2.3/3.0/3.9 at $m=2/3/4$);
LLM judge REPLICATED.}

\end{document}
